\section{The harmonic degree is generically equal to the pencil rank}
\label{deg:sec}

The finite-difference determinant report proves an upper bound for the
harmonic degree and asks when that bound is strict.  For real rational
weights, the answer is particularly clean: it is never strict over the
field of rational functions in the cutoff.  At a specified integer
cutoff it can be strict, but only at finitely many cutoffs.  Positivity
removes these exceptions once enough nodes have been included.

\subsection{An affine symmetric-pencil lemma}

Let $C_1,\ldots,C_t$ be real symmetric $r\times r$ matrices, write
\[
 C(X)=\sum_{j=1}^tX_jC_j,
 \qquad \rho=\operatorname{rank}_{\mathbb R(X)} C(X),
 \qquad q=r-\rho,
\]
and let $D$ be another real symmetric matrix.  The variables $X_j$ are
formal independent variables.  Their introduction makes no assertion
about algebraic independence of special functions.

\begin{lemma}[Compression and the highest degree]
\label{deg:compression}
The polynomial $P(X)=\det(D+C(X))$ has total degree at most $\rho$.
Choose any $v\in\mathbb R^t$ such that $C(v)$ has rank $\rho$, and let
$K$ be a real $r\times q$ matrix whose columns form its kernel.
Then
\begin{equation}\label{deg:isotropic}
 K^{\mathsf T}C_jK=0\qquad(1\le j\le t).
\end{equation}
If $K^{\mathsf T}DK$ is nonsingular, $P$ has total degree exactly $\rho$.
For a generic maximal-rank direction, nonsingularity of this compression
is also necessary.
\end{lemma}
\begin{proof}
In each determinant term, selecting more than $\rho$ columns from
$C(X)$ gives zero, which proves the upper bound.  Put $A=C(v)$ and
choose an invertible real matrix $S=[U\ K]$ so that
\[
 S^{\mathsf T}AS=\begin{pmatrix} A_0&0\\0&0\end{pmatrix},
 \qquad \det A_0\ne0.
\]
For every $j$, the rank of $A+zC_j$ is at most $\rho$ for every real
$z$: every $(\rho+1)$-minor of the whole pencil is identically zero.
For sufficiently small $z$, its upper-left $\rho\times\rho$ block
remains invertible.  The Schur complement of this block must therefore
vanish identically.  Its coefficient of $z$ is exactly
$K^{\mathsf T}C_jK$, proving \eqref{deg:isotropic}.

The coefficient of $z^\rho$ in $\det(D+zA)$ is
\begin{equation}\label{deg:leading-direction}
 \frac{\det A_0}{(\det S)^2}\det(K^{\mathsf T}DK).
\end{equation}
This follows by taking the $zA_0$ entry in each of the first $\rho$
rows and columns after the congruence by $S$; the remaining block is
$K^{\mathsf T}DK$.  Thus a nonzero compression proves degree $\rho$.
Conversely, if the degree-$\rho$ homogeneous part is nonzero, it does
not vanish at a generic real maximal-rank direction, and
\eqref{deg:leading-direction} makes that compression nonsingular.
The conventions for $q=0$ are that $K$ has no columns and its
compressed determinant is one.
\end{proof}

There is a direct symbolic version that avoids expanding the full
$r\times r$ determinant.  Choose a nonsingular principal
$\rho\times\rho$ submatrix of $C(X)$ and, after a simultaneous
permutation of rows and columns, write
\[
 C(X)=\begin{pmatrix} A&B\\B^{\mathsf T}&F\end{pmatrix}.
\]
Over $\mathbb R(X)$, the rank condition gives $F=B^{\mathsf T}A^{-1}B$.
Set
\begin{equation}\label{deg:generic-kernel}
 K(X)=\begin{pmatrix}-A^{-1}B\\I_q\end{pmatrix}.
\end{equation}
Then the highest homogeneous part has the exact expression
\begin{equation}\label{deg:highest}
 [P]_{\rho}(X)=\det A\,
                  \det\bigl(K(X)^{\mathsf T}DK(X)\bigr).
\end{equation}
Although the displayed factors are rational functions, their product is
the indicated homogeneous polynomial.  The expression remains valid
when it vanishes.  A symmetric matrix of rank $\rho$ over a field of
characteristic zero has a nonsingular principal minor of size $\rho$;
this follows, for example, by symmetric elimination with $1\times1$
and $2\times2$ nonsingular pivots.

\subsection{Sharpness for moment matrices, including signed weights}

Let $w_n\in\mathbb R$ be a sequence and put
\begin{equation}\label{deg:moment}
 M_N=\sum_{n=0}^N w_n v(n)v(n)^{\mathsf T},
 \qquad v(n)=(1,n,\ldots,n^{r-1})^{\mathsf T}.
\end{equation}
Suppose that a decomposition
\begin{equation}\label{deg:decomp}
 M_N=D_N+\sum_{j=1}^t h_j(N)C_j
\end{equation}
holds with real $h_j(N)$ and real symmetric matrices $D_N,C_j$,
where the $C_j$ are independent of $N$.  Define
$P_N(X)=\det(D_N+C(X))$ and let $\delta_N$ be its total degree,
with degree $-\infty$ assigned to the zero polynomial.

\begin{theorem}[Only finitely many deficient cutoffs]
\label{deg:main}
Assume there are an integer $N_*\ge0$ and $\epsilon\in\{1,-1\}$ such
that $\epsilon w_n>0$ for every integer $n>N_*$.  Then
\[
 \delta_N\le\rho\quad\text{for all }N,
\]
and at most $r(r-\rho)$ integers $N\ge N_*$ satisfy
$\delta_N<\rho$.
If all nonzero $w_n$ have the same sign and at least $r$ distinct
nodes of positive absolute weight have been included by cutoff $N$,
then $\delta_N=\rho$ at that cutoff.

If in addition $D_N$ has entries rational in $N$, then
\begin{equation}\label{deg:generic-sharp}
 \deg_X\det(D_N+C(X))=\rho
 \qquad\text{over }\mathbb R(N).
\end{equation}
\end{theorem}
\begin{proof}
Choose one real maximal-rank direction and its kernel matrix $K$ as in
Lemma~\ref{deg:compression}.  Such a real direction exists because a
nonzero maximal minor is a nonzero real polynomial.  Equation
\eqref{deg:isotropic} and \eqref{deg:decomp} imply
\begin{equation}\label{deg:gram}
 G_N:=K^{\mathsf T}D_NK
 =K^{\mathsf T}M_NK
 =\sum_{n=0}^N w_n p(n)p(n)^{\mathsf T},
 \qquad p(n)=K^{\mathsf T}v(n).
\end{equation}
The components of $p$ are linearly independent real polynomials of
degree at most $r-1$.

For $N\ge N_*$, the matrix
\[
 \epsilon(G_{N+r}-G_N)
 =\sum_{n=N+1}^{N+r}\epsilon w_n p(n)p(n)^{\mathsf T}
\]
is positive definite.  Indeed, a nonzero linear combination of the
components of $p$ is a nonzero polynomial of degree at most $r-1$;
it cannot vanish at all $r$ displayed nodes.  Along any one of the
$r$ residue classes of cutoffs modulo $r$, the ordered eigenvalues of
$\epsilon G_N$ therefore increase strictly.  Each of the $q=r-\rho$
eigenvalues can equal zero at most once.  Consequently at most $q$
cutoffs in each residue class have singular $G_N$, and at most $rq$
cutoffs in total do.  By Lemma~\ref{deg:compression}, degree deficiency
can occur only at those cutoffs.  This also treats $q=0$, for which
there are no exceptions.

If the weights in the cutoff have one sign and at least $r$ distinct
nonzero nodes, the same polynomial argument makes $G_N$ definite
directly.  Finally, in the rational case $\det G_N$ is a rational
function of $N$.  It is nonzero at some integer cutoff by the preceding
argument, hence is not the zero rational function.  Formula
\eqref{deg:leading-direction} proves \eqref{deg:generic-sharp}.
\end{proof}

\begin{corollary}[Resolution for real rational weights]
\label{deg:rational}
Let $R$ be a nonzero real rational function having no pole at a
nonnegative integer.  Form its finite moment matrices and their
harmonic decomposition as in the finite-difference determinant report.
Then their formal harmonic degree over $\mathbb R(N)$ is exactly the
generic harmonic-pencil rank.  At most finitely many nonnegative
integer cutoffs have smaller degree.  For $R(n)>0$ at all
nonnegative integers, every $N\ge r-1$ has the exact generic degree.
\end{corollary}
\begin{proof}
A nonzero real rational function has constant nonzero sign for all
sufficiently large positive arguments.  The residue construction gives
the decomposition \eqref{deg:decomp} with rational $D_N$.  Apply
Theorem~\ref{deg:main}.  When poles and residues are represented using
nonreal conjugate pairs, replace each conjugate pair of formal harmonic
coordinates by its formal real and imaginary parts.  This invertible
linear change makes the coefficient matrices real and preserves both
generic rank and total degree.
\end{proof}

The real hypothesis in this proof is substantive: it supplies definiteness
of the sufficiently late increments.  No corresponding theorem for an
arbitrary complex rational weight is asserted here.  Likewise, formal
degree does not by itself prove a minimal multiple-polylogarithm depth
under functional relations.

\subsection{An exact test for exceptional cutoffs}

In the generic kernel formula \eqref{deg:generic-kernel}, replace $K$ by
the polynomial matrix
\[
 \widetilde K(X)=
 \begin{pmatrix}-\operatorname{adj}(A)B\\(\det A)I_q\end{pmatrix}.
\]
After clearing denominators in $N$, compute the polynomial
\begin{equation}\label{deg:exception-certificate}
 Q(N,X)=\det\bigl(\widetilde K(X)^{\mathsf T}D_N
                                      \widetilde K(X)\bigr).
\end{equation}
For an admissible cutoff $N_0$, degree is deficient if and only if
$Q(N_0,X)$ is identically zero.  Equivalently, all its coefficients in
$X$ vanish.  Over an exact coefficient field, their greatest common
divisor as polynomials in $N$ gives precisely the exceptional cutoffs
among its admissible nonnegative integer roots.  Denominator zeros must
be excluded or checked from the original finite moment formula.  This
uses a determinant of size $r-\rho$, not an expansion of the original
determinant.  A single fixed maximal-rank direction gives a cheaper
finite candidate list, which the generic test can then filter.

\subsection{A positive example and a genuine signed exception}

Take $r=3$ and first let
\[
 R(x)=\frac1{x+1}+\frac1{x+2}.
\]
There is one harmonic variable $X=H_N^{(1)}(1)$, and its coefficient is
\[
 C=\begin{pmatrix}2&-3&5\\-3&5&-9\\5&-9&17\end{pmatrix},
 \qquad \operatorname{rank}C=2,
 \qquad K=(2,3,1)^{\mathsf T}.
\]
Since $K^{\mathsf T}v(n)=(n+1)(n+2)$, the leading coefficient is
\begin{equation}\label{deg:positive-example}
 [P_N]_2=\frac{(N+1)(N+2)^2(N+3)}2 X^2.
\end{equation}
Here the top-left $2\times2$ minor of $C$ has determinant one.

For the signed weight
\begin{equation}\label{deg:signed-weight}
 R(x)=\frac1{x+1}-\frac7{5(x+2)}
     =\frac{3-2x}{5(x+1)(x+2)},
\end{equation}
the same kernel has compressed moment
\[
 K^{\mathsf T}D_NK
 =-\frac{(N-2)(N+1)(N+2)(N+3)}{10}.
\]
The relevant principal minor of $C$ has determinant $-7/5$, so
\begin{equation}\label{deg:signed-leading}
 [P_N]_2=\frac{7(N-2)(N+1)(N+2)(N+3)}{50}X^2.
\end{equation}
The only exceptional nonnegative cutoff is $N=2$, where the entire
formal polynomial is
\begin{equation}\label{deg:signed-exception}
 P_2(X)=\frac{620X-1137}{500}.
\end{equation}
Substituting the actual $X=H_2^{(1)}(1)=11/6$ gives $-1/1500$, also
equal to the three-node Vandermonde square times
$R(0)R(1)R(2)$.  Thus a fixed-cutoff exception is real, even when the
finite moment determinant itself is nonzero.  It does not contradict
generic sharpness over $\mathbb R(N)$.
